% ==============================================================
% Projektive Strukturtheorie (PST)
% Dokument PST-5-S: Strategie, Zwischenbilanz und offene Fragen
% Version 0.1
% Kompiliert mit: pdfLaTeX
% ==============================================================

\documentclass[11pt,a4paper]{article}

% ---------- Grundpakete ----------
\usepackage[utf8]{inputenc}
\usepackage[T1]{fontenc}
\usepackage[german]{babel}
\usepackage{amsmath,amssymb,amsthm}
\usepackage{geometry}
\usepackage{parskip}
\usepackage{enumitem}
\usepackage{booktabs}
\usepackage{xcolor}
\usepackage{hyperref}

% ---------- Layout ----------
\geometry{margin=2.5cm}
\setlist{itemsep=0.2em, topsep=0.2em}
\hypersetup{
    colorlinks=true,
    linkcolor=blue,
    urlcolor=blue,
    pdftitle={PST-5-S: Strategie und Zwischenbilanz -- Version 0.1},
    pdfauthor={Forschungsprogramm Ontoph\"{a}nomenalismus}
}

% ---------- Theorem-Umgebungen ----------
\theoremstyle{plain}
\newtheorem{theorem}{Theorem}[section]
\newtheorem{proposition}{Proposition}[section]

\theoremstyle{definition}
\newtheorem{definition}{Definition}[section]
\newtheorem{remark}{Bemerkung}[section]
\newtheorem{methodennote}{Methodennote}[section]
\newtheorem{hypothese}{Arbeitshypothese}[section]
\newtheorem{prinzip}{Prinzip}[section]
\newtheorem{ziel}{Forschungsziel}[section]

% Neue Kommandos
\newcommand{\OPP}{\mathcal{M}_{\mathrm{OPP}}}
\newcommand{\deff}{d_{\mathrm{eff}}}
\newcommand{\Proj}{\mathcal{P}}

% ---------- Titel ----------
\title{
  \textbf{Projektive Strukturtheorie (PST)}\\[0.4em]
  \textbf{Dokument PST-5-S: Strategie, Zwischenbilanz
  und offene Fragen}\\[0.4em]
  {\large Forschungsstrategie, priorisierte Agenda und
  Selbstverst\"{a}ndnis der PST}\\[0.5em]
  {\normalsize Version~0.1}
}
\author{
  \textbf{Forschungsprogramm Ontoph\"{a}nomenalismus}\\
  \small Siehe TRANSPARENCY f\"{u}r methodologische Urheberschaft
  und KI-Einsatz.
}
\date{Juli 2026}

\begin{document}

\maketitle

\begin{abstract}
Dieses Dokument zieht eine Zwischenbilanz der PST, entwickelt die
Forschungsstrategie f\"{u}r die n\"{a}chsten Phasen und benennt die
priorisierten offenen Fragen.

Das Selbstverst\"{a}ndnis der PST: Sie ist kein abgeschlossenes System,
sondern ein offenes Forschungsprogramm im Sinne von Imre Lakatos. Der
harte Kern --- die Rekonstruktionskette
Existenz $\rightarrow$ OPP $\rightarrow \mathcal{P} \rightarrow$
Raumzeit $\rightarrow$ Physik --- ist durch das No-Go Ergebnis und die
Diffusion-Maps-Befunde gest\"{a}rkt worden. Der Schutzg\"{u}rtel aus
Hilfs\-hypothesen wird laufend angepasst.

Das Fernziel bleibt:
\[
  \boxed{\text{Existenz} \;\Longrightarrow\; \text{Universum.}}
\]

Es gelten die methodologischen Konventionen gem\"{a}\ss{} MKO.
\end{abstract}

\tableofcontents
\newpage

%==============================================================
\section{Selbstverst\"{a}ndnis: Die PST als Forschungsprogramm}
%==============================================================

\begin{methodennote}[PST im Sinne von Lakatos]
Die PST versteht sich als wissenschaftliches Forschungsprogramm im Sinne
von Imre Lakatos. Ein Forschungsprogramm besteht aus:
\begin{itemize}
  \item \textbf{Hartem Kern:} Unver\"{a}nderliche Grundhypothesen die
        das Programm definieren.
  \item \textbf{Schutzg\"{u}rtel:} Hilfshypothesen die angepasst werden
        k\"{o}nnen wenn Evidenz dagegen spricht.
  \item \textbf{Heuristik:} Methodische Regeln die bestimmen, wie das
        Programm weiterentwickelt wird.
\end{itemize}
Ein Programm ist \emph{progressiv} wenn neue Versionen mehr voraussagen
als alte. Es ist \emph{degenerativ} wenn es nur ad hoc auf Gegenbeispiele
reagiert. Die PST ist bisher progressiv: Das No-Go Ergebnis f\"{u}hrte
zu einer tieferen Theorie, nicht zu einer ad hoc Rettung.
\end{methodennote}

\begin{definition}[Harter Kern der PST]
Der harte Kern der PST --- die Annahmen die nicht aufgegeben werden
k\"{o}nnen ohne das Programm aufzugeben:
\begin{enumerate}
  \item Der OPP existiert als hochdimensionaler, pr\"{a}-metrischer
        Zustandsraum.
  \item Raumzeit, Physik und Bestimmtheit emergieren durch Projektion.
  \item Es gibt einen universellen Projektionsoperator $\Proj^\ast$.
  \item Das Prinzip der minimalen Ontologie gilt.
\end{enumerate}
\end{definition}

\begin{definition}[Schutzg\"{u}rtel der PST]
Anpassbare Hilfshypothesen:
\begin{enumerate}
  \item Diffusion Maps sind $\Proj^\ast$ (ersetzbar durch anderen Operator).
  \item $t \approx 9$--$10$ ist der physikalische Skalenparameter
        (ersetzbar durch pr\"{a}zisere Herleitung).
  \item $I_0 \leftrightarrow \hbar$ (ersetzbar durch pr\"{a}zisere
        Operationalisierung).
  \item Mutual Information als Relationsst\"{a}rke (ersetzbar durch
        andere Gr\"{o}\ss{}e).
\end{enumerate}
\end{definition}

%==============================================================
\section{Zwischenbilanz: Was die PST bisher geleistet hat}
%==============================================================

\begin{remark}[Positive Leistungen der PST]
Die PST hat bisher geleistet:
\begin{enumerate}
  \item \textbf{Pr\"{a}zisierung der Frage:} Von \glqq{}Hat der OPP 4D?\grqq{}
        zu \glqq{}Welche Projektionen erzeugen 4D?\grqq{} --- ein
        fundamentaler konzeptueller Fortschritt.
  \item \textbf{Das No-Go Ergebnis:} Der rohe OPP hat keine intrinsische
        Vierdimensionalit\"{a}t. Das schr\"{a}nkt den Theorieraum ein.
  \item \textbf{Das $\deff \approx 4$ Signal:} Diffusion Maps mit
        $t \approx 9$--$10$ liefern reproduzierbar $\deff = 3.98 \pm 0.36$
        ohne eingebaute Vierdimensionalit\"{a}t. Das ist der bisher
        st\"{a}rkste empirische Befund.
  \item \textbf{Das Symmetriebrechungsprinzip:} Projektion = Symmetriebrechung
        ist ein neues konzeptuelles Prinzip mit Vorhersagekraft.
  \item \textbf{Methodische Reinigung:} Rangneutralit\"{a}t als Kriterium;
        Unterscheidung lokaler und globaler Verfahren; Selbstpr\"{u}fung
        als Pflicht.
  \item \textbf{Verbindung zu OP und DPR:} Die PST hat die philosophischen
        Konzepte von DPR (Projektionsoperatoren, RdE) mathematisch
        ausgearbeitet.
\end{enumerate}
\end{remark}

\begin{remark}[Grenzen der bisherigen PST]
Die PST hat bisher nicht geleistet:
\begin{enumerate}
  \item Den universellen Operator $\Proj^\ast$ hergeleitet.
  \item Die Lorentz-Signatur $(+,-,-,-)$ aus Projektion emergiert.
  \item Die Eichgruppen $SU(3) \times SU(2) \times U(1)$ erzeugt.
  \item Quantisierung als Projektionsartefakt bewiesen.
  \item Eine einzige Naturkonstante aus OPP-Eigenschaften abgeleitet.
  \item Dynamik aus dem statischen OPP erkl\"{a}rt.
\end{enumerate}
Diese Liste ist ehrlich und wird in v1.0 der PST-Dokumente
als Forschungsagenda dienen.
\end{remark}

%==============================================================
\section{Die Forschungsstrategie: Top-Down und Bottom-Up}
%==============================================================

\begin{prinzip}[Doppelte Strategie]
Die PST verfolgt gleichzeitig zwei komplementäre Ans\"{a}tze:
\begin{itemize}
  \item \textbf{Top-Down:} Ausgangspunkt ist die bekannte Physik.
        Gesucht werden Projektionsoperatoren die Lorentz-Invarianz,
        Eichgruppen und Quantisierung reproduzieren.
  \item \textbf{Bottom-Up:} Ausgangspunkt ist der OPP. Gesucht werden
        emergente Strukturen ohne physikalische Vorannahmen.
\end{itemize}
Beide Ans\"{a}tze sind notwendig. Top-Down liefert Randbedingungen;
Bottom-Up liefert neue, unerwartete Strukturen.
\end{prinzip}

\begin{remark}[Warum Bottom-Up allein scheitert]
Reines Bottom-Up --- man beginnt mit dem OPP und hofft auf Physik ---
ist gescheitert (No-Go Ergebnis). Der OPP allein legt keine physikalische
Struktur nahe. Man braucht die Randbedingungen der Physik als Leitfaden.
\end{remark}

\begin{remark}[Warum Top-Down allein gef\"{a}hrlich ist]
Reines Top-Down --- man sucht Operatoren die die bekannte Physik
reproduzieren --- l\"{a}uft Gefahr zirkul\"{a}r zu werden: Man findet,
was man sucht. Das Rangneutralit\"{a}tsprinzip (PST-4-G) ist die
Schutzma\ss{}nahme gegen diesen Fehler.
\end{remark}

%==============================================================
\section{Priorisierte Forschungsagenda}
%==============================================================

\begin{ziel}[Priorit\"{a}t 1: $\Proj^\ast$ aus Prinzipien]
\textbf{Frage:} Kann der universelle Projektionsoperator aus dem Prinzip
der minimalen Ontologie und den Randbedingungen der Physik eindeutig
hergeleitet werden?

\textbf{Warum Priorit\"{a}t 1:} Ohne $\Proj^\ast$ bleibt die PST
deskriptiv. Mit $\Proj^\ast$ wird sie prädiktiv --- sie kann Physik
vorhersagen, nicht nur rekonstruieren.

\textbf{N\"{a}chste Schritte:}
\begin{itemize}
  \item Charakterisierung der Klasse von Projektionsoperatoren die
        $\deff = 4$ liefern.
  \item Pr\"{u}fung welcher Operator die Lorentz-Signatur reproduziert.
  \item Kategorientheoretische Formulierung von $\Proj^\ast$ als Funktor.
\end{itemize}
\end{ziel}

\begin{ziel}[Priorit\"{a}t 2: Lorentz-Signatur]
\textbf{Frage:} Wie emergiert $(+,-,-,-)$ aus einer signaturfreien
Gram-Matrix?

\textbf{Warum Priorit\"{a}t 2:} $\deff \approx 4$ ist ein Signal f\"{u}r
vier Dimensionen --- aber noch nicht f\"{u}r die Lorentz-Signatur.
Der Unterschied zwischen euklidischer 4D ($SO(4)$) und Minkowski-4D
($SO(1,3)$) ist physikalisch entscheidend.

\textbf{Kandidatenansatz:} Imagin\"{a}re Eigenwerte bei komplexen
Gram-Matrizen k\"{o}nnten die Signatur erzeugen.
\end{ziel}

\begin{ziel}[Priorit\"{a}t 3: Mathematische Fundierung]
\textbf{Frage:} Welcher mathematische Formalismus beschreibt OPP,
Attraktoren und Projektionsoperatoren am pr\"{a}zisesten?

\textbf{Kandidaten:} Kategorientheorie, Operatoralgebren, Homotopietheorie,
spektrale Tripel (Connes).

\textbf{Warum Priorit\"{a}t 3:} Ohne pr\"{a}zise Mathematik bleiben
alle Aussagen \"{u}ber $\Proj^\ast$ vage. Die mathematische Fundierung
ist Voraussetzung f\"{u}r alle anderen Ziele.
\end{ziel}

\begin{ziel}[Priorit\"{a}t 4: Eichgruppen]
\textbf{Frage:} Wie emergieren $SU(3) \times SU(2) \times U(1)$ als
Symmetriegruppen der physikalischen Projektion?

\textbf{Warum Priorit\"{a}t 4:} Die Eichgruppen des Standardmodells
sind die pr\"{a}ziseste bekannte Beschreibung der Materiestruktur.
Ihre Ableitung aus $\Proj^\ast$ w\"{a}re ein Durchbruch.
\end{ziel}

\begin{ziel}[Priorit\"{a}t 5: Quantisierung und Naturkonstanten]
\textbf{Frage:} Ist $I_0 \leftrightarrow \hbar$ pr\"{a}zisierbar?
K\"{o}nnen Naturkonstanten aus OPP-Eigenschaften abgeleitet werden?

\textbf{Warum Priorit\"{a}t 5:} Naturkonstanten sind die h\"{a}rtesten
Pr\"{u}fsteine jeder Fundamentaltheorie. Eine einzige korrekt abgeleitete
Konstante w\"{a}re ein starker Beleg f\"{u}r die PST.
\end{ziel}

%==============================================================
\section{Kriterien f\"{u}r Fortschritt und Scheitern}
%==============================================================

\begin{methodennote}[Wann ist die PST erfolgreich?]
Die PST ist erfolgreich wenn:
\begin{enumerate}
  \item $\Proj^\ast$ eindeutig aus Prinzipien hergeleitet wird.
  \item $\deff = 4$ mathematisch notwendig folgt (nicht nur numerisch).
  \item Lorentz-Signatur aus $\Proj^\ast$ emergiert.
  \item Mindestens eine Naturkonstante korrekt abgeleitet wird.
\end{enumerate}
\end{methodennote}

\begin{methodennote}[Wann scheitert die PST?]
Die PST scheitert wenn:
\begin{enumerate}
  \item Es sich zeigt, dass $\deff \approx 4$ ein Artefakt der
        Diffusion-Maps-Implementierung ist (kontrollierte Tests laufen).
  \item Kein rangneutraler Operator $\Proj^\ast$ existiert der alle
        Randbedingungen erf\"{u}llt.
  \item Die Lorentz-Signatur grundsätzlich nicht aus einer reellen,
        positiv semidefiniten Gram-Matrix emergieren kann.
\end{enumerate}
Jedes dieser Scheitern-Kriterien ist pr\"{u}fbar. Die PST ist falsifizierbar
im Sinne von Popper.
\end{methodennote}

%==============================================================
\section{Das Fernziel und seine Bedeutung}
%==============================================================

\begin{remark}[Was das Fernziel bedeutet]
Das Fernziel der PST lautet:
\[
  \text{Existenz} \;\Longrightarrow\; \text{Universum.}
\]
Dies bedeutet: Aus der einzigen gesicherten Tatsache --- dass
\"{u}berhaupt etwas existiert --- sollen alle Eigenschaften des
physikalischen Universums mathematisch notwendig folgen.

Das ist das ambiti\"{o}seste Ziel das ein Forschungsprogramm sich stellen
kann. Es ist m\"{o}glicherweise nicht vollst\"{a}ndig erreichbar. Aber
jeder Schritt in diese Richtung --- jede Eigenschaft des Universums die
aus minimalen Annahmen folgt --- ist ein echter wissenschaftlicher
Fortschritt.
\end{remark}

\begin{remark}[Was die PST \emph{nicht} beansprucht]
Die PST beansprucht nicht:
\begin{itemize}
  \item Eine fertige Theorie zu sein.
  \item Alle physikalischen Fragen zu beantworten.
  \item Die einzig m\"{o}gliche mathematische Ausarbeitung des OP zu sein.
  \item Konkurrenztheorien (String-Theorie, Loop-Quanten-Gravitation,
        kausal-dynamische Triangulierung) zu widerlegen.
\end{itemize}
Die PST beansprucht: Die konsequenteste mathematische Ausarbeitung des
philosophischen Programms des OP zu sein, die mit dem Prinzip der
minimalen Ontologie konsistent ist.
\end{remark}

%==============================================================
\section{Verh\"{a}ltnis zu anderen Ans\"{a}tzen}
%==============================================================

\begin{remark}[PST und String-Theorie]
Die String-Theorie beginnt mit einem Stringobjekt --- das ist bereits
eine geometrische Annahme. Die PST beginnt mit dem OPP ohne geometrische
Vorannahmen. Der Unterschied ist fundamental: String-Theorie setzt
Geometrie voraus, PST leitet sie ab. Ob beide Programme letztlich
\"{a}quivalent sind, ist eine offene Frage.
\end{remark}

\begin{remark}[PST und Loop-Quanten-Gravitation (LQG)]
LQG diskretisiert die Raumzeit --- Raum besteht aus endlichen
Quantenvolumina. Die PST macht keine solche Diskretisierungsannahme:
Der OPP ist kontinuierlich; Diskretisierung emergiert durch die endliche
Aufl\"{o}sung des Projektionsoperators. Dies k\"{o}nnte eine tiefere
Erkl\"{a}rung f\"{u}r das LQG-Bild liefern.
\end{remark}

\begin{remark}[PST und nicht-kommutative Geometrie (Connes)]
Alain Connes beschreibt Raumzeit \"{u}ber spektrale Tripel $(A, H, D)$
--- eine Algebra $A$, ein Hilbert-Raum $H$ und einen Dirac-Operator $D$.
Die PST k\"{o}nnte diese Struktur als einen spezifischen Projektionsoperator
interpretieren. Connes' Ansatz ist m\"{o}glicherweise eine mathematisch
ausgearbeitete Version des PST-Projektionsprinzips.
\end{remark}

%==============================================================
\section{Schnittstellen}
%==============================================================

\begin{itemize}
  \item \textbf{PST-1-F bis PST-4-G:} PST-5-S ist die Synthese aller
        vorangegangenen PST-Dokumente. Es zieht die Bilanz und gibt
        die Richtung vor.
  \item \textbf{PST-6-R:} Die priorisierte Forschungsagenda in Sektion~4
        ist das Arbeitsprogramm f\"{u}r PST-6-R.
  \item \textbf{OP:} Das Fernziel der PST ist die Einl\"{o}sung des
        philosophischen Versprechens des OP: Wenn das Denkfeld der
        Urgrund ist, muss die Physik daraus folgen.
  \item \textbf{MGO:} Priorit\"{a}t~3 (mathematische Fundierung) ist
        eine direkte Aufgabe f\"{u}r MGO~v1.1.
  \item \textbf{AR (Adversarial Review):} Die Scheitern-Kriterien in
        Sektion~5 sind direkt als AR-Testfragen verwendbar.
\end{itemize}

\end{document}

